\documentclass[12pt, twoside]{article}
%\documentclass[12pt, twoside]{article}
\usepackage[letterpaper, margin=1in, headsep=0.2in]{geometry}
\setlength{\headheight}{0.6in}
%\usepackage[english]{babel}
\usepackage[utf8]{inputenc}
\usepackage{microtype}
\usepackage{amsmath}
\usepackage{amssymb}
%\usepackage{amsfonts}
\usepackage[nomessages]{fp} %\FPeval{\var-name}{2*sin(pi/6)}
\usepackage{siunitx} %units in math. eg 20\milli\meter
\usepackage{yhmath} % for arcs, overparenth command
\usepackage{tikz} %graphics
\usetikzlibrary{quotes, angles, arrows, arrows.meta}
\usepackage{pgfplots}
\usepgfplotslibrary{fillbetween}
\pgfplotsset{compat=1.16}
\usepackage{graphicx} %consider setting \graphicspath{{images/}}
\usepackage{parskip} %no paragraph indent
\usepackage{enumitem}
\usepackage{multicol}
\usepackage{venndiagram}

\usepackage{fancyhdr}
\pagestyle{fancy}
\fancyhf{}
\renewcommand{\headrulewidth}{0pt} % disable the underline of the header
\raggedbottom
\hfuzz=2mm %suppresses overfull box warnings

\usepackage{hyperref}

\fancyhead[LE]{\thepage}
\fancyhead[RO]{Markscheme}
\fancyhead[LO]{La Scuola d'Italia / Dr.\ Huson / IB Math 11 \\* 29 September 2026}

\begin{document}
\subsubsection*{1.6 Classwork: Geometric Growth and Decay --- Markscheme}

\begin{enumerate}[itemsep=0.3cm]

%--- Problem 1: Common ratio, u_n, u_10, first term over 1 000 000 [6 marks] ---
\item[1.] [Maximum mark: 6]

The sequence $4, 12, 36, 108, \ldots$, so $u_{1} = 4$.

\medskip

\textbf{(a)} $\dfrac{12}{4} = 3, \quad \dfrac{36}{12} = 3, \quad
\dfrac{108}{36} = 3$ \hfill \textbf{M1}

$r = 3$ \hfill \textbf{A1}

\emph{M1 for at least two of the three ratios correctly calculated.
A1 for $r = 3$.}

\emph{A common difference in place of a ratio (e.g.\ $12 - 4 = 8$)
earns no marks.}

\medskip

\textbf{(b)} $u_{n} = 4 \times 3^{\,n - 1}$ \hfill \textbf{A1}

\emph{Accept any equivalent form, e.g.\ $\dfrac{4}{3} \times 3^{n}$.}

\emph{$u_{n} = 4 \times 3^{n}$ or $12^{\,n-1}$ earns no mark.}

\emph{FT: award A1 for $4 \times r^{\,n-1}$ with the candidate's own
$r$ from (a).}

\medskip

\textbf{(c)} $u_{10} = 4 \times 3^{9} = 78\,732$ \hfill \textbf{A1}

\emph{FT: award A1 for the value obtained by substituting $n = 10$
into the candidate's own expression from (b).}

\medskip

\textbf{(d)} $4 \times 3^{\,n - 1} > 1\,000\,000$ \hfill \textbf{M1}

$u_{12} = 708\,588 \ (< 1\,000\,000)$ \quad and \quad
$u_{13} = 2\,125\,764 \ (> 1\,000\,000)$

$n = 13$ \hfill \textbf{A1}

\emph{Accept solving the inequality on the GDC or with logarithms:
$n - 1 > 11.31\ldots$, so $n = 13$.}

\emph{Accept a table search: full marks for showing the adjacent rows
$u_{12}$ and $u_{13}$ and stating $n = 13$. A single row
($u_{13} = 2\,125\,764$, so $n = 13$) still earns full marks, with the
feedback note that the adjacent row $u_{12}$ is what justifies
``smallest''. See markscheme-spec §18.}

\emph{$n = 12.31\ldots$ is not an acceptable final answer --- $n$ must
be a whole number. Rounding it down to $n = 12$ earns the M1 only.}

\emph{FT: award A1 for the smallest integer $n$ correctly obtained
from the candidate's own expression in (b).}

\newpage

%--- Problem 2: Population growth 4% and car depreciation 15% [7 marks] ---
\item[2.] [Maximum mark: 7]

Population: $12\,000$ at the start of 2020, growing $4\%$ per
year. Car: \$$24\,000$ at the start of 2020, losing $15\%$ per year.

\medskip

\textbf{(a)} $r = 1.04$ \hfill \textbf{A1}

\emph{Accept $1 + 0.04$ or $\dfrac{104}{100}$.}

\emph{$r = 0.04$ or $r = 4$ earns no mark.}

\medskip

\textbf{(b)} 2030 is $10$ years after 2020:
$12\,000 \times 1.04^{10}$ \hfill \textbf{M1}

$= 17\,762.9\ldots \approx 17\,763$ \hfill \textbf{A1}

\emph{Accept $17\,800$ (3 s.f.), $17\,762$, $17\,762.9$ or any value
that rounds to $17\,760$.}

\emph{An exponent of $9$ (giving $17\,079.7\ldots$) or $11$ (giving
$18\,473.4\ldots$) earns the M1 only.}

\emph{FT: award M1A1 for $12\,000 \times r^{10}$ evaluated with the
candidate's own $r$ from (a), provided $r > 1$.}

\medskip

\textbf{(c)} $r = 0.85$ \qquad
$24\,000 \times 0.85^{5}$ \hfill \textbf{M1}

$=$ \$$10\,648.93 \approx$ \$$10\,600$ \hfill \textbf{A1}

\emph{Accept \$$10\,600$ (3 s.f.), \$$10\,649$ or \$$10\,648.93$. M1 is
for a multiplier of $0.85$ applied five times, however written.}

\emph{A multiplier of $0.15$ (giving \$$1.82$) or $1.15$ earns no
marks. An exponent of $4$ (giving \$$12\,528.15$) earns the M1 only.}

\medskip

\textbf{(d)} $24\,000 \times 0.85^{\,n} < 5000$ \hfill \textbf{M1}

After $9$ years: \$$5558.81$ \ ($> 5000$); \quad after $10$ years:
\$$4724.99$ \ ($< 5000$)

$10$ years \hfill \textbf{A1}

\emph{Accept solving the inequality on the GDC or with logarithms:
$n > 9.65\ldots$, so $10$ years.}

\emph{Accept a table search: full marks for showing both adjacent
years and stating $10$. A single row (after $10$ years the value is
\$$4724.99$, so $10$) still earns full marks, with the feedback note
that the adjacent year $9$ is what justifies ``first''. See
markscheme-spec §18.}

\emph{$9.65$ is not an acceptable final answer --- the answer must be
a whole number of years. Rounding down to $9$ earns the M1 only.}

\emph{FT: award A1 for the smallest whole number of years correctly
obtained from the candidate's own multiplier in (c), provided it is
between $0$ and $1$.}

\end{enumerate}
\end{document}
